\subsection{Cantor--Lebesgue Singular Function}

\begin{definition} \label{def:cantor-lebesgue-singular-function-source} \hlt{Cantor--Lebesgue Singular Function}\\
Let $P\subseteq[0,1]$ be the Cantor set. For $x=(0.c_1c_2\cdots)_3\in P$ with $c_k\in\{0,2\}$, set $\varphi(x)\defeq(0.(c_1/2)(c_2/2)\cdots)_2$.\\
If $(a,b)$ is a component of $[0,1]\setminus P$, extend by $\varphi(x)\defeq\varphi(a)=\varphi(b)$ for $x\in(a,b)$.\\
The resulting $\varphi:[0,1]\to[0,1]$ is the \hlt{Cantor--Lebesgue singular function}.
\end{definition}

\begin{proposition} \label{prop:cantor-singular-function-properties} \hlt{Properties of the Singular Function}\\
Let $P\subseteq[0,1]$ be the Cantor set and let $\varphi:[0,1]\to[0,1]$ be the Cantor--Lebesgue singular function.\\
Then $\varphi$ is continuous, nondecreasing, onto, nonconstant, and $\varphi'=0$ on $[0,1]\setminus P$.
\begin{proof} On $P$, the digit map $\tau((0.c_1c_2\cdots)_3)\defeq(0.(c_1/2)(c_2/2)\cdots)_2$ is nondecreasing and onto $[0,1]$.\\
The only duplications are endpoint pairs of removed middle-third intervals, where the binary values agree; constant extension across those intervals is therefore nondecreasing with no jumps, hence continuous.\\
Each component of $[0,1]\setminus P$ is an interval of constancy, so $\varphi'=0$ there. Since $\varphi([0,1])=[0,1]$, $\varphi$ is nonconstant.\\
\hlt{Structural Commentary}: Variation is confined to $P$; the range is $[0,1]$.\\
\hlt{Key Insight}: Ternary $0/2$ digits are re-read as binary digits.
\end{proof}
\end{proposition}

\begin{theorem} \label{thm:lebesgue-measurable-non-borel-set-exists} \hlt{Lebesgue Measurable Non-Borel Sets Exist}\\
There exist $A\in\mathcal{L}^1\setminus\mathcal{B}(\R)$ and a homeomorphism $\psi:[0,1]\to[0,2]$ such that $\psi(A)\notin\mathcal{L}^1$.
\begin{proof} Let $\psi(x)\defeq x+\varphi(x)$, where $\varphi$ is the Cantor--Lebesgue singular function. Since $\varphi$ is continuous and nondecreasing, $\psi$ is continuous, strictly increasing, maps $[0,1]$ onto $[0,2]$, and hence is a homeomorphism.\\
Let $P$ be the Cantor set and $U\defeq[0,1]\setminus P$. On each component $I$ of $U$, $\varphi$ is constant, so $m(\psi(I))=m(I)$; the intervals $\psi(I)$ are disjoint. Hence $m(\psi(U))=m(U)=1$.\\
Since $[0,2]=\psi(P)\sqcup\psi(U)$ and $\psi(P)$ is compact, $m(\psi(P))=1$. Choose $W\subseteq\psi(P)$ with $W\notin\mathcal{L}^1$, and set $A\defeq\psi^{-1}(W)\subseteq P$. Since $m(P)=0$, $A\in\mathcal{L}^1$. If $A\in\mathcal{B}(\R)$, then $W=\psi(A)$ is Borel in $[0,2]$, hence Lebesgue measurable, contradiction. Thus $A\in\mathcal{L}^1\setminus\mathcal{B}(\R)$ and $\psi(A)=W\notin\mathcal{L}^1$.\\
\hlt{Structural Commentary}: Completion by null sets creates Lebesgue sets invisible to the Borel hierarchy.\\
\hlt{Key Insight}: A null Cantor set can be homeomorphically stretched to positive measure.
\end{proof}
\end{theorem}
